\documentclass[12pt]{jlreq}
\usepackage{amsmath, amssymb}

\newcommand{\C}{\mathbb{C}}
\newcommand{\R}{\mathbb{R}}
\newcommand{\Z}{\mathbb{Z}}
\newcommand{\N}{\mathbb{N}}
\newcommand{\F}{\mathbb{F}}
\newcommand{\Prob}{\mathbb{P}}
\newcommand{\Exp}{\mathbb{E}}
\newcommand{\dd}{\,d}
\newcommand{\Ker}{\operatorname{Ker}}
\newcommand{\Impart}{\operatorname{Im}}
\newcommand{\Hom}{\operatorname{Hom}}

\begin{document}

\(x>0\) に対し \(\Gamma(x)=\displaystyle\int_0^\infty e^{-t}t^{x-1}\dd t\) とおく。

\begin{enumerate}
  \item 正則関数 \(f(w)=e^{-tw}w^{x-1}\) の領域
  \[
    D_\varepsilon=\left\{w\in\C\ \middle|\ \operatorname{Re}w>0,\ \operatorname{Im}w>0,\ \varepsilon<|w|<\frac1\varepsilon\right\},\qquad \varepsilon>0,
  \]
  の境界での積分を考えることにより
  \[
    \int_0^\infty \sin t\cdot t^{x-1}\dd t
    =\Gamma(x)\sin\left(\frac{\pi x}{2}\right),\qquad 0<x<1,
  \]
  を示せ。
  \item
  \[
    \int_0^\infty \frac{\sin t}{t^{3/2}}\dd t=\sqrt{2\pi}
  \]
  を示せ。必要であれば，\(\Gamma(1/2)=\sqrt{\pi}\) を用いてもよい。
\end{enumerate}

\end{document}